\documentclass[11pt]{article}
\usepackage[margin=1in]{geometry}
\usepackage[T1]{fontenc}
\usepackage[utf8]{inputenc}
\usepackage{amsmath,amssymb,amsthm}
\usepackage{enumitem}

\title{ICPC World Finals 2022\\S. Bridging the Gap}
\author{}
\date{}

\begin{document}
\maketitle

\section*{Problem Summary}

There are $n$ walkers with crossing times $t_1,\dots,t_n$ and a bridge that can hold at most $c$ walkers
at once. A forward crossing carries the torch to the far side; unless it is the last forward crossing, some
walker must later bring the torch back. We must minimize the total time until everyone is across.

\section*{Key Observations}

\begin{itemize}[leftmargin=*]
    \item Sort the walkers by time:
    \[
    t_1 \le t_2 \le \cdots \le t_n.
    \]
    In an optimal solution, every backward crossing is made by exactly one walker: if several walkers came
    back together, keeping only the fastest of them is never worse.
    \item Before the last trip, a forward crossing consists of some number of currently slowest walkers that
    are meant to stay on the far side, plus some number of fastest walkers that will later return one by one.
    Using any non-fast walker as a future returner is dominated by using a faster one.
    \item This suggests a dynamic program over:
    \begin{itemize}[leftmargin=*]
        \item how many slow walkers have already been moved permanently;
        \item how many future one-person return trips have already been ``paid for'' by sending fast helpers
        across earlier.
    \end{itemize}
    \item Besides mixed trips that move slow walkers, we may also take a \emph{fast-only package}: send
    $h+1$ fastest walkers across, planning for $h$ of them to return later. This produces $h$ future return
    credits and costs
    \[
    t_{h+1} + \sum_{i=1}^{h} t_i.
    \]
    These package costs can be precomputed with an unbounded knapsack.
\end{itemize}

\section*{Algorithm}

\begin{enumerate}[leftmargin=*]
    \item Sort the crossing times increasingly.
    \item Precompute \texttt{credit\_cost[x]} = minimum extra cost of fast-only packages that create exactly
    $x$ future return credits.
    \item Let \texttt{dp[done][credit]} be the minimum cost after the \texttt{done} slowest walkers have been
    moved permanently, with \texttt{credit} net prepaid returns available.
    \item From state \texttt{(done, credit)}, choose how many fast helpers accompany the next block of slow
    walkers. If \texttt{helpers = h}, then:
    \begin{itemize}[leftmargin=*]
        \item the crossing moves $c-h$ of the remaining slow walkers permanently;
        \item its time contribution is the slowest moved walker plus $\sum_{i=1}^{h} t_i$, because those
        $h$ fast helpers will later return individually;
        \item the net number of credits changes by $h-1$, since every non-final forward crossing needs one
        return trip.
    \end{itemize}
    \item The DP only keeps the credit range that can still matter, which reduces the number of reachable
    states to $O(n^2/c)$.
    \item After all slow walkers have been transferred, add enough fast-only packages to settle any remaining
    credit deficit and take the best final cost.
\end{enumerate}

\section*{Correctness Proof}

We prove that the algorithm returns the correct answer.

\paragraph{Lemma 1.}
There exists an optimal solution in which every backward crossing consists of a single walker.

\paragraph{Proof.}
Suppose some backward crossing contains several walkers. Replacing it by only the fastest walker from
that group still returns the torch to the near side and cannot increase the time of that crossing. The other
walkers can simply remain on the far side. Repeating this transformation yields an optimal solution with
only single-person returns. \qed

\paragraph{Lemma 2.}
There exists an optimal solution in which every forward crossing is either:
\begin{itemize}[leftmargin=*]
    \item a group of the currently slowest remaining walkers, possibly together with some of the globally
    fastest walkers that will later return; or
    \item a fast-only package.
\end{itemize}

\paragraph{Proof.}
By Lemma 1, future returners always cross back alone. If a forward trip contains a walker who will later
return, replacing that walker by a faster not-yet-used walker can only decrease future return costs and
cannot increase the forward trip time. Thus temporary helpers may be assumed to be chosen from the
globally fastest walkers.

The walkers intended to stay permanently on the far side may then be taken in decreasing order of speed:
if two permanent walkers are transferred in the opposite order, swapping them does not worsen any trip,
because only the slowest walker in each forward group matters. Hence some optimal solution has exactly
the stated structure. \qed

\paragraph{Lemma 3.}
Every schedule of the form described in Lemma 2 is represented by a unique path in the dynamic
program, and every DP path corresponds to a valid schedule of that form.

\paragraph{Proof.}
The state records precisely the two quantities needed to continue optimally: how many slow walkers have
already been transferred permanently, and how many prepaid one-person returns are still available. Each
transition chooses the number of fast helpers on the next forward trip, which uniquely determines how
many slow walkers are moved, what time is added, and how the credit balance changes. Conversely, any
schedule of the structured form yields exactly those state changes. Therefore the DP and the structured
schedules are in one-to-one correspondence. \qed

\paragraph{Theorem.}
The algorithm outputs the minimum total crossing time.

\paragraph{Proof.}
By Lemmas 1 and 2, some optimal solution has the structured form modeled by the DP. By Lemma 3, the
DP explores exactly all such solutions and computes their exact total cost, including the cheapest possible
fast-only packages needed to settle the final credit balance. Therefore the minimum value produced by the
DP is the true optimum. \qed

\section*{Complexity Analysis}

The pruned state space has size $O(n^2/c)$, and each state has $O(c)$ outgoing transitions, so the main
DP runs in $O(n^2)$ time. The fast-only package preprocessing is also $O(n^2)$ in the worst case. The
memory usage is $O(n^2/c)$.

\section*{Implementation Notes}

\begin{itemize}[leftmargin=*]
    \item The code stores only the credit range that can still matter for a given \texttt{done}, which is why
    each DP row has a different width.
    \item If $c \ge n$, everybody can cross together immediately and the answer is simply the slowest
    walker's time.
\end{itemize}

\end{document}
